\documentclass[11pt,a4paper]{article}
\usepackage[T1]{fontenc}
\usepackage[utf8]{inputenc}
\usepackage[main=italian,provide=*]{babel}
\usepackage{amsmath,amssymb}
\usepackage{geometry}
\geometry{margin=2.3cm,top=2.5cm,bottom=2.7cm}
\usepackage{booktabs}
\usepackage[table]{xcolor}
\usepackage{tcolorbox}
\tcbuselibrary{skins,breakable}
\usepackage{titlesec}
\usepackage{enumitem}
\usepackage{seqsplit}
\usepackage{needspace}
\usepackage{fancyhdr}
\usepackage{hyperref}
\hypersetup{colorlinks=true,linkcolor=Sepia,citecolor=Sepia,urlcolor=Sepia}

\definecolor{inkblue}{HTML}{1B3A5C}
\definecolor{lightblue}{HTML}{EAF1F8}
\definecolor{accent}{HTML}{B0562A}
\definecolor{softgray}{HTML}{6B6B6B}
\definecolor{okgreen}{HTML}{2E6B3E}
\definecolor{okgreenbg}{HTML}{EAF5EC}
\definecolor{warnamber}{HTML}{9C6B12}
\definecolor{warnbg}{HTML}{FBF1E0}
\definecolor{advisorpurple}{HTML}{5B3A7A}
\definecolor{advisorbg}{HTML}{F1ECF6}
\definecolor{notebar}{HTML}{9C6B12}

\titleformat{\section}
  {\normalfont\Large\bfseries\color{inkblue}}
  {\thesection}{0.6em}{}[{\color{inkblue!40}\titlerule}]
\titlespacing{\section}{0pt}{0.8em}{0.4em}
\titleformat{\subsection}
  {\normalfont\large\bfseries\color{accent}}
  {}{0em}{}
\titlespacing{\subsection}{0pt}{0.5em}{0.2em}

\pagestyle{fancy}
\fancyhf{}
\renewcommand{\headrulewidth}{0.4pt}
\fancyhead[R]{\small\color{softgray}\thepage}
\fancyfoot[C]{\small\color{softgray}Tesi triennale, Samuele Schiavi --- Relatore: Prof.\ Alessandro Chiesa}

\newtcolorbox{keyeq}[1][]{
  colback=lightblue, colframe=inkblue!70, boxrule=0.6pt,
  arc=2pt, left=8pt, right=8pt, top=4pt, bottom=4pt, #1
}
\newtcolorbox{panelbox}[1]{
  colback=white, colframe=softgray!50, boxrule=0.5pt,
  arc=2pt, left=7pt, right=7pt, top=3pt, bottom=3pt,
  fonttitle=\bfseries\color{inkblue}, title={#1}
}
\newtcolorbox{okbox}[1]{
  colback=okgreenbg, colframe=okgreen!60, boxrule=0.6pt,
  arc=2pt, left=7pt, right=7pt, top=3pt, bottom=3pt,
  fonttitle=\bfseries\color{okgreen}, title={#1}
}
\newtcolorbox{warnbox}[1]{
  colback=warnbg, colframe=warnamber!70, boxrule=0.6pt,
  arc=2pt, left=7pt, right=7pt, top=3pt, bottom=3pt,
  fonttitle=\bfseries\color{warnamber}, title={#1}
}
\newtcolorbox{advisorbox}[1]{
  colback=advisorbg, colframe=advisorpurple!60, boxrule=0.6pt,
  arc=2pt, left=7pt, right=7pt, top=4pt, bottom=4pt,
  fonttitle=\bfseries\color{advisorpurple}, title={#1}
}
\newtcolorbox{editnote}{
  colback=white, colframe=white, boxrule=0pt,
  borderline west={2.2pt}{0pt}{notebar},
  arc=0pt, left=10pt, right=4pt, top=2pt, bottom=2pt,
  fontupper=\itshape\small\color{softgray!90!black}
}
\fancyhead[L]{\small\color{softgray}Circuito compatto per il passo di Trotter}

\title{\vspace{-1.5em}\color{inkblue}Circuito compatto per il passo di Trotter\\[0.1em]
\Large\normalfont\color{softgray}Teoria, verifica, e perché non si usa (per ora)}
\author{Note di lavoro --- Tesi triennale, Samuele Schiavi\\ Relatore: Prof.\ Alessandro Chiesa}
\date{}

\begin{document}
\sloppy
\maketitle
\thispagestyle{fancy}
\vspace{-2em}

Nota di riferimento su un'analisi collaterale fatta durante il lavoro
sul dimero: il Prof.\ Chiesa ha chiesto se il circuito del passo di
Trotter potesse essere reso più compatto (esempio: con 2 CNOT invece
degli 8 gate a due qubit attuali). Questo file raccoglie la teoria
usata, la verifica numerica fatta (con due metodi indipendenti), il
risultato, e --- soprattutto --- perché quel risultato \textbf{non va
usato ingenuamente} quando si passerà a discutere di rumore (sistema
aperto) e di ricerca del ground state. Pensato per essere caricato come
riferimento all'inizio di quella discussione.

\section{Il problema}

Il circuito attuale per un singolo passo di Trotter usa \textbf{8 gate
a due qubit}: $1\times R_{XX}+1\times R_{YY}+6\times\text{CNOT}$ (i 6
CNOT vengono dai tre blocchi $R_{ZZ}$ scritti esplicitamente come
CNOT--$R_z$--CNOT: uno per lo scambio, due per il termine DM).
Profondità 15.

La domanda del relatore: si può fare con meno gate, ad esempio con 2
CNOT?

\section{La teoria: quanti CNOT servono per un'unitaria a 2 qubit qualsiasi}

\subsection*{Il risultato generale (letteratura)}

Ogni operatore unitario a 2 qubit $U\in SU(4)$ si può scrivere, a meno
di trasformazioni locali a un qubit (che non costano CNOT), nella forma
canonica di Cartan:
\begin{keyeq}
\centering
$U_{\rm local} = \exp\!\left[-\dfrac{i}{2}\big(c_1\,X\!\otimes\!X + c_2\,Y\!\otimes\!Y + c_3\,Z\!\otimes\!Z\big)\right]$,
\quad $0\le c_3\le c_2\le c_1\le\dfrac\pi2$
\end{keyeq}
I tre numeri $(c_1,c_2,c_3)$ --- le \textbf{coordinate di Weyl} ---
classificano completamente la "non-località" del gate,
indipendentemente da come è vestito con rotazioni a un qubit. Il
numero minimo di CNOT necessari a realizzare $U$ dipende \textbf{solo}
da queste coordinate:

\begin{center}
\renewcommand{\arraystretch}{1.25}
\begin{tabular}{@{}lc@{}}
\toprule
\rowcolor{inkblue!10}
\textbf{Condizione su $(c_1,c_2,c_3)$} & \textbf{CNOT minimi}\\
\midrule
$c_1=c_2=c_3=0$ & 0 (gate locale, nessun entanglement)\\
$(c_1,c_2,c_3)=(\pi/4,0,0)$ & 1 (classe di equivalenza locale del CNOT)\\
$c_3=0$ (ma non nel caso sopra) & 2\\
$c_3\neq0$ & 3 (il caso generico)\\
\bottomrule
\end{tabular}
\end{center}

\subsection*{Bibliografia --- verificata con fonti indipendenti (non solo plausibile)}

Ogni citazione qui sotto è stata controllata contro fonti terze
(pagine ufficiali degli editori, bibliografie di altri lavori) prima
di essere inclusa:

\begin{itemize}[leftmargin=1.3em,itemsep=0.4em]
\item \textbf{F.\ Vatan, C.\ Williams}, \emph{Optimal Quantum Circuits
for General Two-Qubit Gates}, Phys.\ Rev.\ A \textbf{69}, 032315
(2004). Dimostra che 3 CNOT bastano sempre e sono necessari nel caso
generico, con costruzione esplicita indipendente. \textbf{Riferimento
primario consigliato per questo teorema}: nessuna correzione nota in
letteratura.
\item \textbf{G.\ Vidal, C.\ M.\ Dawson}, \emph{A Universal Quantum
Circuit for Two-Qubit Transformations With Three CNOT Gates}, Phys.\
Rev.\ A \textbf{69}, 010301 (2004). Stesso risultato, costruzione
diversa.
\begin{editnote}
Attenzione: esiste un Comment pubblicato --- M.\ W.\ Coffey, R.\
Deiotte, T.\ Semi, \emph{Phys.\ Rev.\ A} \textbf{77}, 066301 (2008) ---
che segnala errori nella costruzione esplicita presentata in questo
paper. Non è stato possibile verificare il dettaglio dell'errore
(testo a pagamento), ma il teorema principale (3 CNOT necessari e
sufficienti nel caso generico) resta lo standard accettato: citato
senza riserve in decine di lavori successivi fino al 2024--2025, ed è
comunque corroborato indipendentemente dal punto precedente
(Vatan--Williams) e dalla verifica numerica fatta qui sotto con il
codice di Qiskit.
\end{editnote}
\item \textbf{Y.\ Makhlin}, \emph{Nonlocal Properties of Two-Qubit
Gates and Mixed States and Optimization of Quantum Computations},
Quantum Inf.\ Process.\ \textbf{1}, 243 (2002) --- introduce gli
invarianti locali (equivalenti alle coordinate di Weyl).
\item \textbf{J.\ Zhang, J.\ Vala, S.\ Sastry, K.\ B.\ Whaley},
\emph{Geometric theory of nonlocal two-qubit operations}, Phys.\ Rev.\
A \textbf{67}, 042313 (2003) --- lega gli invarianti di Makhlin alla
decomposizione di Cartan e alla "camera di Weyl".
\end{itemize}

\subsection*{Perché conta $c_3$ e non $c_1,c_2$}

Intuitivamente: $c_1$ e $c_2$ si possono sempre "assorbire" in parte
con trasformazioni locali astute; è la terza coordinata, quella più
piccola nella convenzione ordinata $c_1\ge c_2\ge c_3$, a essere
l'ostacolo residuo che nessuna rotazione locale può eliminare. Un CNOT
(o CZ, equivalente a meno di locali) ha $(c_1,c_2,c_3)=(\pi/4,0,0)$:
$c_3=0$ esattamente, il caso più semplice possibile per un gate
davvero entanglante. Un'operazione con $c_3\neq0$ non può quindi mai
essere raggiunta con un solo CNOT, né con due (che restano confinati a
$c_3=0$), e servono 3.

\needspace{10\baselineskip}
\section{Verifica del metodo su casi noti}

Prima di applicare il criterio al passo di Trotter, è stato validato
su quattro casi di cui il risultato è già certo, usando
\texttt{qiskit.synthesis.TwoQubitWeylDecomposition}:

\begin{center}
\renewcommand{\arraystretch}{1.25}
\begin{tabular}{@{}lccc@{}}
\toprule
\rowcolor{inkblue!10}
\textbf{Gate} & \textbf{Coord.\ di Weyl $(a,b,c)$} & \textbf{CNOT minimi (atteso)} & \textbf{Verificato}\\
\midrule
identità & $(0,0,0)$ & 0 & \checkmark\\
CNOT & $(\pi/4,0,0)$ & 1 & \checkmark\\
iSWAP & $(\pi/4,\pi/4,0)$ & 2 & \checkmark\\
SWAP & $(\pi/4,\pi/4,\pi/4)$ & 3 & \checkmark\\
\bottomrule
\end{tabular}
\end{center}

\section{Applicazione al passo di Trotter del dimero}

Con $b/J=-0.18$, $D/J=1$ (il punto di lavoro adottato) e $\tau$
variabile:
$$U_{\rm step}(\tau) = e^{-iH_2\tau}\,e^{-iH_1\tau}$$

\begin{center}
\renewcommand{\arraystretch}{1.25}
\begin{tabular}{@{}cccc@{}}
\toprule
\rowcolor{inkblue!10}
$\tau$ & Weyl $(a,b,c)$ & CNOT minimi & \\
\midrule
0.10 & $(0.035,\,0.035,\,-0.025)$ & 3 & \\
0.50 & $(0.173,\,0.173,\,-0.125)$ & 3 & \\
1.00 & $(0.319,\,0.319,\,-0.250)$ & 3 & \\
1.50 & $(0.391,\,0.391,\,-0.375)$ & 3 & \\
\bottomrule
\end{tabular}
\end{center}

La terza coordinata $c\neq0$ per \textbf{ogni} $\tau\neq0$ testato (e
per costruzione, per ogni $\tau\neq0$ in generale, dato che
$c\propto-\tau/4$ in questo caso --- proporzionale al termine DM).
Verificato anche nel caso $D=0$ (senza termine DM): $c\neq0$ comunque,
quindi nemmeno l'assenza del DM riporta al caso a 2 CNOT --- lo scambio
isotropo da solo già impone $c\neq0$.

\subsection*{Seconda verifica, con un metodo Qiskit indipendente dal primo}

Oltre a \texttt{TwoQubitWeylDecomposition}, il risultato è stato
ricontrollato con \texttt{TwoQubitBasisDecomposer} (un algoritmo di
sintesi diverso, non solo una lettura delle coordinate), forzando
esplicitamente un numero fissato di CNOT e misurando la fedeltà
effettivamente raggiunta rispetto a $U_{\rm step}(\tau=1)$:

\begin{center}
\renewcommand{\arraystretch}{1.25}
\begin{tabular}{@{}cc@{}}
\toprule
\rowcolor{inkblue!10}
\textbf{CNOT forzati} & \textbf{Fedeltà raggiunta}\\
\midrule
0 & $0.764$\\
1 & $0.676$\\
2 & $0.939$\\
3 & $1.000$\\
\bottomrule
\end{tabular}
\end{center}

Fedeltà esatta \textbf{solo} a partire da 3 CNOT --- mai con 2, per
quanta ottimizzazione numerica si tenti. Il metodo
\texttt{num\_basis\_gates()} di Qiskit, che implementa autonomamente
questa stessa ricerca, restituisce \textbf{3} in modo indipendente.

\begin{okbox}{Conclusione sulla domanda originale}
2 CNOT sono matematicamente impossibili per questo passo, confermato
da due algoritmi diversi. Il minimo raggiungibile è 3.
\end{okbox}

\needspace{10\baselineskip}
\section{Dal minimo teorico (3) al circuito realizzato: due strade diverse}

\subsection*{Comprimere il singolo passo}

Il circuito attuale usa 8 gate a due qubit per passo; il minimo è 3.
Il circuito compatto è stato ottenuto in tre passaggi distinti, ed è
importante tenerli separati per capire cosa produce davvero:

\begin{enumerate}[leftmargin=1.4em,itemsep=0.25em]
\item \textbf{Calcolo classico preliminare}: la matrice $4\times4$
esatta del passo, $U_{\rm step}=e^{-iH_2\tau}e^{-iH_1\tau}$, calcolata
con \texttt{scipy.linalg.expm} --- non un circuito, pura algebra
lineare classica.
\item \textbf{La matrice passata come blocco numerico opaco}:
\texttt{qc.unitary(U\_step, [0,1])} --- a questo punto Qiskit ha in
mano solo 16 numeri complessi, non sa che provengono da
un'Hamiltoniana con un termine di scambio e un termine DM.
\item \textbf{Sintesi via transpiler}:
\texttt{transpile(qc, basis\_gates=['cx','rz','sx','x'],
optimization\_level=3)}, che internamente usa la decomposizione di
Cartan/KAK descritta in sez.~2 per ricostruire \emph{qualunque} matrice
$4\times4$ con il minor numero di CNOT possibile.
\end{enumerate}

\begin{editnote}
Il punto cruciale è nel passaggio 2--3: l'ottimizzatore non sa cosa sia
$H_1$, $H_2$, lo scambio isotropo o il termine DM --- vede solo "questa
matrice $4\times4$, riproducila con meno gate possibile". È pura
ottimizzazione numerica sulla matrice finale, \emph{indifferente a come
quella matrice si è generata}. Per questo il circuito compatto che ne
esce ha angoli $R_z$ senza alcun significato fisico leggibile (non si
legge più "questo pezzo è lo scambio, questo è il DM") --- non è un
dettaglio implementativo, è la conseguenza diretta del fatto che
l'informazione sulla struttura fisica è stata scartata già al
passaggio 2, quando l'Hamiltoniana è stata ridotta a una matrice di
numeri senza etichette. Il file \texttt{circuito\_compatto.png} mostra
il risultato di questi tre passaggi, con fedeltà $1.0$ rispetto al
passo esatto.
\end{editnote}

\subsection*{Comprimere l'intero circuito a $N$ passi --- il punto concettualmente più importante}

Qui emerge un fatto che a prima vista sembra sorprendente: comprimendo
l'\textbf{intero} prodotto
$U_{\rm Trotter}(t,N)=\big[e^{-iH_2\tau}e^{-iH_1\tau}\big]^N$ (non passo
per passo, ma il prodotto finale), il risultato resta \textbf{sempre
un'unica unitaria a 2 qubit} --- e quindi si comprime anch'esso a
\textbf{3 CNOT, indipendentemente da $N$}:

\begin{center}
\renewcommand{\arraystretch}{1.25}
\begin{tabular}{@{}ccccc@{}}
\toprule
\rowcolor{inkblue!10}
$N$ & gate 2q (esplicito) & gate 2q (compresso) & risparmio & \\
\midrule
1 & 8 & 3 & 62\% & \\
2 & 16 & 3 & 81\% & \\
5 & 40 & 3 & 92\% & \\
20 & 160 & 3 & 98\% & \\
100 & 800 & 3 & 99.6\% & \\
\bottomrule
\end{tabular}
\end{center}

Verificato: la fedeltà del circuito compresso rispetto a
$U_{\rm Trotter}(t,N)$ (non rispetto all'evoluzione esatta) resta
$1.0$ per ogni $N$ --- cioè il circuito compresso realizza
\textbf{esattamente} la stessa unitaria approssimata di Trotter, solo
con meno gate. L'errore di Trotter (rispetto all'evoluzione vera) è
identico prima e dopo la compressione, a qualunque $N$: comprimere
\textbf{non cambia l'accuratezza}, cambia solo quanti gate fisici
servono per realizzarla.

\needspace{12\baselineskip}
\section{Perché questa scorciatoia non regge per il lavoro con il rumore (e va dichiarata come tale)}

Questo è il punto da avere chiaro \textbf{prima} di passare alla Parte
2.

\begin{warnbox}{Il compromesso $N$-vs-rumore sparisce}
Tutta l'idea della Parte 2 è studiare come l'errore \emph{algoritmico}
di Trotter (che scala come $\sim1/N$, e quindi vuole $N$ grande) e
l'errore \emph{fisico} dei gate sotto un modello di rumore (che scala
con il numero di gate, e quindi vuole $N$ piccolo) si bilancino fra
loro, per trovare un compromesso ottimale. Se si usa il circuito
compresso, il costo in gate resta fisso a 3 \textbf{qualunque sia $N$}:
il compromesso a cui la Parte 2 è interessata smette di esistere nel
grafico. Si misurerebbe il rumore su un'unitaria diversa dall'oggetto
di studio.
\end{warnbox}

\subsection*{La compressione richiede un calcolo classico pregresso, quindi non è "vero" Trotter}

Lo stesso meccanismo descritto per il singolo passo (calcolo classico
della matrice, poi sintesi) si applica identico all'intero prodotto:
per ottenere il circuito a 3 CNOT bisogna \textbf{prima} calcolare
classicamente $U_{\rm Trotter}(t,N)$ come matrice $4\times4$ (con
\texttt{scipy}), e solo \emph{dopo} comprimerla in gate. Su un sistema
più grande di 2 qubit --- cioè il caso per cui Trotter è stato
inventato --- questo calcolo classico preliminare è esattamente quello
che si vuole evitare (matrice $2^n\times2^n$, intrattabile per $n$
grande). La compressione funziona qui solo perché il dimero è un
giocattolo a 2 qubit comodo da diagonalizzare a mano: è una
scorciatoia specifica di questo caso di test, non un metodo
generalizzabile a sistemi più grandi.

\subsection*{Dove la forma compatta avrebbe senso, con l'avvertenza giusta}

Se in futuro si vuole discutere "qual è il rumore minimo possibile per
realizzare \emph{questa specifica} evoluzione a un tempo $t$ fissato"
(una domanda a sé, legittima ma diversa), il circuito compresso è lo
strumento giusto. Ma va tenuto esplicitamente \textbf{separato} dal
test "vero" di Trotter con rumore, che deve usare il circuito esplicito
(o comunque uno il cui conteggio di gate scali con $N$), altrimenti si
confondono due esperimenti concettualmente diversi.

\needspace{12\baselineskip}
\section{Riepilogo per la discussione futura}

Quando si affronterà il rumore su un sistema quantistico aperto
(generico su X, Y, Z) e la ricerca del ground state, tenere presente:

\begin{enumerate}[leftmargin=1.4em,itemsep=0.3em]
\item \textbf{Il conteggio di gate del circuito esplicito scala
linearmente con $N$} (8 gate a due qubit per passo, $8N$ in totale) ---
è quello giusto da usare per studiare il compromesso Trotter/rumore,
perché riflette il costo fisico reale che si avrebbe anche su un
sistema più grande.
\item \textbf{Il circuito compresso a 3 CNOT è una scorciatoia valida
solo per la verifica puntuale su questo sistema-giocattolo a 2 qubit}
--- utile per sapere "qual è il limite fisico minimo", ma da non usare
come sostituto del circuito di Trotter quando si introduce il rumore,
pena perdere il fenomeno che si vuole studiare.
\item Se il lavoro sul ground state passerà per metodi variazionali
(VQE) o evoluzione in tempo immaginario, la stessa distinzione tornerà
utile: un circuito "compresso/ottimizzato numericamente" che riproduce
un risultato target non è la stessa cosa di un circuito costruito con
una struttura fisica interpretabile (ansatz), anche quando le due
unitarie coincidono esattamente --- la prima non generalizza, non è
ispezionabile termine a termine, e non dice nulla sul comportamento
sotto rumore realistico a scala più grande.
\end{enumerate}

\needspace{15\baselineskip}
\section*{Appendice: comandi Qiskit usati per la verifica}
\addcontentsline{toc}{section}{Appendice}

\begin{panelbox}{Codice}
\small
\begin{verbatim}
from qiskit.synthesis import TwoQubitWeylDecomposition, \
    TwoQubitBasisDecomposer
from qiskit.circuit.library import CXGate
from qiskit.quantum_info import Operator, process_fidelity
from qiskit import QuantumCircuit, transpile

# metodo 1: coordinate di Weyl di un'unitaria U (matrice numpy 4x4)
d = TwoQubitWeylDecomposition(Operator(U))
a, b, c = d.a, d.b, d.c   # c=0 -> 2 CNOT bastano; c!=0 -> servono 3

# metodo 2 (indipendente): sintesi forzata a N CNOT, verifica fedelta'
decomposer = TwoQubitBasisDecomposer(CXGate())
for n in [0, 1, 2, 3]:
    qc = decomposer(Operator(U), _num_basis_uses=n)
    print(n, process_fidelity(Operator(qc), Operator(U)))
print("minimo autonomo:", decomposer.num_basis_gates(Operator(U)))

# compressione automatica a CNOT minimi (per generare il circuito finale)
qc = QuantumCircuit(2)
qc.unitary(U, [0, 1])
compresso = transpile(qc, basis_gates=['cx', 'rz', 'sx', 'x'],
                       optimization_level=3)
\end{verbatim}
\end{panelbox}

\end{document}
